\documentclass[12pt, twoside]{article}
%\documentclass[12pt, twoside]{article}
\usepackage[letterpaper, margin=1in, headsep=0.2in]{geometry}
\setlength{\headheight}{0.6in}
%\usepackage[english]{babel}
\usepackage[utf8]{inputenc}
\usepackage{microtype}
\usepackage{amsmath}
\usepackage{amssymb}
%\usepackage{amsfonts}
\usepackage[nomessages]{fp} %\FPeval{\var-name}{2*sin(pi/6)}
\usepackage{siunitx} %units in math. eg 20\milli\meter
\usepackage{yhmath} % for arcs, overparenth command
\usepackage{tikz} %graphics
\usetikzlibrary{quotes, angles, arrows, arrows.meta}
\usepackage{pgfplots}
\usepgfplotslibrary{fillbetween}
\pgfplotsset{compat=1.16}
\usepackage{graphicx} %consider setting \graphicspath{{images/}}
\usepackage{parskip} %no paragraph indent
\usepackage{enumitem}
\usepackage{multicol}
\usepackage{venndiagram}

\usepackage{fancyhdr}
\pagestyle{fancy}
\fancyhf{}
\renewcommand{\headrulewidth}{0pt} % disable the underline of the header
\raggedbottom
\hfuzz=2mm %suppresses overfull box warnings

\usepackage{hyperref}

\fancyhead[LO]{La Scuola d'Italia / Dr.\ Huson / IB Math 11 \\* 5 October 2026}
\fancyhead[RO]{\fbox{\begin{minipage}{6cm}\footnotesize Nickname:\\[0.5cm]\end{minipage}}}
\fancyhead[LE]{\fbox{\begin{minipage}{6cm}\footnotesize Nickname:\\[0.5cm]\end{minipage}}}
\fancyfoot[C]{\thepage}

\begin{document}
\subsubsection*{1.8 Classwork: Compound interest and inflation}

\begin{enumerate}[itemsep=0.15cm]

%--- Problem 1: Compound interest compounded annually; value, interest, first year over $8000 [6 marks] ---
%--- Source: Original (freshly authored, AA SL 1.4) ---
\item[1.]

If \$$P$ is invested at $r\%$ per year, compounded $k$ times per year,
the value after $n$ years is
$$FV = P\left(1 + \frac{r}{100k}\right)^{kn}.$$

Sofia invests \$$5000$ in an account that pays $4\%$ per year compounded
annually.

\begin{enumerate}[itemsep=0.1cm]
  \item Write down the number that the value is multiplied by each year.
  \item Find the value of Sofia's investment at the end of $6$ years.
  \item Find the total interest Sofia has earned at the end of $6$ years.
  \item Find the number of complete years after which the value of the
  investment first exceeds \$$8000$.
\end{enumerate}

\begin{tikzpicture}
  \draw (-0.6,0) rectangle (14.6,13.4);
  \foreach \y in {1,2,...,13} \draw[dotted] (0,\y)--(14.1,\y);
\end{tikzpicture}

\newpage

%--- Problem 2: Monthly compounding; real value after inflation; real value of cash [7 marks] ---
%--- Source: Original (freshly authored, AA SL 1.4), modelled on 25M.2.AHL.TZ2.4 ---
\item[2.]

Daniel invests \$$12\,000$ in an account that pays $5.2\%$ per year
compounded \textbf{monthly}.

\begin{enumerate}[itemsep=0.1cm]
  \item Find the value of Daniel's investment at the end of $10$ years.
\end{enumerate}

\textbf{Inflation} is the rise in prices over time, so the same amount of
money buys less in the future. If inflation is $i\%$ per year, then \$$A$
in $n$ years' time buys the same as
$$\frac{A}{\left(1 + \dfrac{i}{100}\right)^{n}}$$
today. This is called the \textbf{real value} of \$$A$.

Over the $10$ years, inflation is $2.8\%$ per year.

\begin{enumerate}[itemsep=0.1cm]
  \item[(b)] Find the real value of Daniel's investment at the end of
  $10$ years.
  \item[(c)] Hence find the real increase in the value of Daniel's
  investment.
  \item[(d)] Daniel's sister keeps \$$12\,000$ in cash at home for the
  same $10$ years. Find the real value of her cash at the end of the
  $10$ years.
\end{enumerate}

\begin{tikzpicture}
  \draw (-0.6,0) rectangle (14.6,11.2);
  \foreach \y in {1,2,...,11} \draw[dotted] (0,\y)--(14.1,\y);
\end{tikzpicture}

\end{enumerate}
\end{document}
